% Reorganized from the frozen October 15, 2026 supplement.
\section{Broadcast uncertainty and queue area}
\label{app:broadcast}

\paragraph{Precise claim.}
Let $0<\varepsilon\le1/16$, $0<\delta\le1/16$,
$0<\eta<\varepsilon/2$. The reference arm has mean
$1/4+\varepsilon/2$. All candidates have mean $p=1/4$ in environment zero;
environment $j$ raises only candidate $j$ to $p+2\varepsilon$. Let $M_t$ be
the candidate count, $a_t$ the new-candidate count, and
\[
 q_- =\left\lfloor\frac{\operatorname{kl}(1/4,\delta)}
                         {\operatorname{kl}(1/4,1/4+2\varepsilon)}\right\rfloor,
 \quad Q_t=(Q_{t-1}+q_-a_t-1)_+,\quad Q_0=0,
\]
\[
 \mathcal B_T=\sum_{t\le T}(M_t-1)_+1\{Q_t>0\}.
\]
Every policy reliable on these environments obeys
\[
 E_0\mathcal U_T(\eta)\ge(3/4-\delta)\mathcal B_T,
 \qquad
 P_0\{\mathcal U_T(\eta)\ge\mathcal B_T/4\}\ge1/4
\]
when $\mathcal B_T>0$.

\begin{proof}
The denominator defining $q_-$ is at most
$\operatorname{kl}(1/4,3/8)\le1/15$, whereas its numerator is at least
$\frac14\log4+\frac34\log(4/5)\ge7/48>2/15$. Thus $q_-\ge2$.
The recursion gives, by induction,
\[
 Q_t=\max_{0\le u\le t}\left\{
       q_-\sum_{r=u+1}^ta_r-(t-u)\right\}.
\]
If $Q_t>0$, there is a suffix in which total candidate demand exceeds all
physical observation slots. Its arms had no observations before that suffix.
Consequently at least one deterministic candidate label $j$ satisfies
$E_0N_j(t)<q_-$. The choice is made from expected counts, not from a realized
history. For any other candidate $i$, membership of $i$ in $\mathcal P_t$ is erroneous
under environment $j$. Its probability there is at most $\delta$.
The history KL is strictly below $\operatorname{kl}(1/4,\delta)$, so data
processing and monotonicity of binary KL force $P_0(i\in \mathcal P_t)<1/4$.
Since all candidates are genuinely good under environment zero,
$P_0(i\in \mathcal N_t)\le\delta$. It follows that
$P_0(i\in \mathcal V_t)\ge3/4-\delta$. These candidates all have boundary distance
$\varepsilon/2>\eta$. Summing at deterministic busy rounds proves the
expectation bound, entirely under the one environment zero.

Let $Y$ count candidate unresolved area only at busy rounds with $M_t\ge2$.
Then $Y\le\mathcal U_T$, $0\le Y\le2\mathcal B_T$, and
$E_0Y\ge11\mathcal B_T/16$. If
$p=P_0(Y\ge\mathcal B_T/4)$, then
$11\mathcal B_T/16\le\mathcal B_T/4+(7\mathcal B_T/4)p$, so $p\ge1/4$.
\end{proof}

\paragraph{Queue area.}
For an integer $q\ge2$, suppose $m$ candidates arrive by time $s$ and put
$D=s+qm-1$. Then
\[
 \sum_{t\le D}(M_t-1)_+1\{Q_t(q)>0\}
 \ge\frac{q-1}{2}m(m-1).
\]
To prove this, remove all subsequent arrivals, which can only decrease both
the queue and its weights. The $qm$ service units finish by $D$. At the
$n$th actual service, at least $\lceil n/q\rceil$ candidates have arrived,
so service rounds weighted by $M_t-1$ sum to at least
$\sum_{n=1}^{qm}(\lceil n/q\rceil-1)=qm(m-1)/2$.
The post-service queue is zero at the last service round of each busy period.
Their cumulative arrival counts are distinct integers in $\{1,\ldots,m\}$,
so removing their weights costs at most $m(m-1)/2$. This proves the claim.

In particular, if along a fixed infinite schedule a policy reliable in every
environment has a deterministic area envelope $a_T=o(T)$, then
$K_t=o(\sqrt t)$. Otherwise along $s_n$ there are at least $c\sqrt{s_n}$
candidates. Choose the first $m_n=\lfloor c\sqrt{s_n}\rfloor$; eventually
$q_-m_n\le s_n$. The queue bound and the preceding probability statement at
$2s_n$ give $P_0(\mathcal U_{2s_n}\ge c's_n)\ge1/4$, contradicting
$\delta\le1/16$ and the proposed envelope.
